\clearpage
\begin{trapbox}
\textbf{C092\quad 常见误区、条件边界与最小反例}
\begin{center}
% Same coordinate scales/ranges. H0:x²-y²/3=1; H1:x²/2-y²/9=1.
% Both pass through P=(2,3), slopes 2 and 3 respectively.
\begin{tikzpicture}[x=.65cm,y=.65cm,font=\small]
\foreach \offset/\case/\title in {0/0/共焦：正交,10/1/非共焦：不正交} {
\begin{scope}[shift={(\offset,0)}]
\node at (1.3,5.8) {\title};
\draw[blue!75!black,thick,domain=10:105,samples=80,variable=\t]
 plot ({4*cos(\t)},{sqrt(12)*sin(\t)});
\ifnum\case=0
\draw[red!75!black,thick,domain=.2:5,samples=80,variable=\t]
 plot ({sqrt(1+\t*\t/3)},\t);
\draw[orange!80!black,dashed,thick] (1,1)--(3,5);
\draw (2.358,2.821)--(2.537,3.179)--(2.179,3.358);
\node at (.5,5.1) {$90^\circ$};
\node at (1.4,-.65) {$k_Ek_H=-1$};
\else
\draw[red!75!black,thick,domain=.2:5,samples=80,variable=\t]
 plot ({sqrt(2+2*\t*\t/9)},\t);
\draw[orange!80!black,dashed,thick] (1.3,.9)--(2.7,5.1);
\draw[gray] ({2+.7*cos(71.56505118)},{3+.7*sin(71.56505118)})
 arc[start angle=71.56505118,end angle=153.43494882,radius=.7];
\node at (.5,5.1) {$\approx81.9^\circ$};
\node at (1.4,-.65) {$k_Ek_{H_1}=-\tfrac32$};
\fi
\draw[teal!85!black,dashed,thick] (-1,4.5)--(4.5,1.75);
\fill (2,3) circle (2.5pt);
\node[anchor=east] at (1.1,2.35) {$P(2,3)$};
\end{scope}}
\end{tikzpicture}


{\small 共同中心、焦轴和交点都不能替代共焦；两图保持相同坐标比例。}
\end{center}
\begin{enumerate}
\item \textbf{把焦角误当切线交角。}本定理 $90^\circ$ 指两切线夹角，
不是 $\angle F_1PF_2$。例1中 $\cos\angle F_1PF_2=3/5\ne0$。
与 C090 区分：C090 研究焦角，C092 研究切线正交。

\item \textbf{“共轴”冒充“共焦”。}判断共焦必须计算焦距：
椭圆 $c_E^2=a^2-b^2$，双曲线 $c_H^2=A^2+B^2$。
同中心、同轴标准式下等价于 $a^2-A^2=b^2+B^2$。
反例：$x^2/16+y^2/12=1$ 和 $x^2/2-y^2/9=1$ 均过 $(2,3)$，
却有 $c_E^2=4$、$c_H^2=11$，切线斜率积为 $-3/2$。

\item \textbf{混淆切线与法线。}由 $r_1+r_2$、$r_1-r_2$ 求到的
$\mathbf u_1\pm\mathbf u_2$ 是\emph{法线方向}；
不能直接把它们称作切向量。两法线垂直最终才推出两切线垂直。

\item \textbf{以为必须求交点坐标。}最快证明只需要“单位向量的和、差点积为零”。
坐标法是标准验证方法，不是几何本质。

\item \textbf{在双曲线左支误用带符号距离差。}定义始终是
$|PF_1-PF_2|=2A$；左、右支只会改变差的\emph{符号}，
不会改变双曲线在该点的法线\emph{方向直线}。

\item \textbf{不检查斜率是否存在。}一般垂直直线可能有竖直切线，
不能无条件写 $k_1k_2=-1$。
在本题 $a>c>A>0$ 的共焦条件下，公共点满足 $x_0y_0\ne0$，
所以斜率确实全部存在；高阶推广时优先用法向量点积。

\item \textbf{把不满足参数条件的情况也代入。}
先检查 $a>c>A>0$，再用 $b^2=a^2-c^2$、$B^2=c^2-A^2$ 求另两个参数。
若 $a=c$ 或 $A=c$，就会出现 $b=0$ 或 $B=0$，标准方程的分母为零；
若 $A=0$，双曲线方程同样不能使用；若 $c=0$，两个焦点重合。
这些情况均不在本结论范围内。

\end{enumerate}
\end{trapbox}
\clearpage
\begin{trapbox}
\textbf{易错提醒（续）：参数与条件检查}
\begin{enumerate}[start=8]
\item \textbf{误以为“有公共点”就有平方关系。}
$E:x^2/16+y^2/12=1$ 与 $H:x^2/2-y^2/9=1$
虽然在 $(2,3)$ 相交，但 $a^2-A^2=16-2=14$，
$b^2+B^2=12+9=21$，\textbf{不相等}，也不共焦。
正确逻辑：\emph{共焦 $\Rightarrow$ 平方关系；公共点本身不能推出平方关系。}

\item \textbf{把 $B$ 误当双曲线纵轴截距。}
$H:x^2/A^2-y^2/B^2=1$ 的 $B$ 叫\emph{虚半轴}，
双曲线本身不与 $y$ 轴相交；$B$ 不等于 $y$ 轴截距。

\item \textbf{把平方关系误用于错位坐标轴的两条曲线。}
$a^2-A^2=b^2+B^2$ 作为共焦的充要判据，需要中心相同、
实轴/长轴方向同为 $x$ 轴；若存在平移、旋转等，还必须核对焦点位置。

\item \textbf{以为只有一个交点或切点。}满足 $a>c>A>0$ 的共焦椭圆、双曲线
有\emph{四个不同实交点}，且每个都是直角相交；正交不等于相切。
\end{enumerate}

\medskip
\textbf{术语补充（不影响使用结论）：}
有些资料把 $b=0$、$B=0$ 等使曲线形状发生改变的情况称为“退化”，
把本题满足 $a>c>A>0$、$b,B>0$ 的椭圆与双曲线称为“非退化”。
其中 $c=0$ 只表示两个焦点重合，椭圆可以成为圆，不宜把所有不符合本题条件的情况都叫作“退化”。
理解本题只需核对参数和焦点位置，无须先记这个术语。

\end{trapbox}
